% Portable theorem guard: define only what the host preamble is missing, so this
% file drops into any of our manuscripts and re-running it is a no-op.
\makeatletter
\@ifundefined{theorem}{\newtheorem{theorem}{Theorem}}{}
\@ifundefined{lemma}{\newtheorem{lemma}[theorem]{Lemma}}{}
\@ifundefined{proposition}{\newtheorem{proposition}[theorem]{Proposition}}{}
\@ifundefined{corollary}{\newtheorem{corollary}[theorem]{Corollary}}{}
\@ifundefined{definition}{\newtheorem{definition}[theorem]{Definition}}{}
\@ifundefined{remark}{\newtheorem{remark}[theorem]{Remark}}{}
\makeatother
% BT1163 quotient Boolean module theorem.

\begin{theorem}[Projected Boolean/Clifford quotient module]
\label{thm:projected-boolean-quotient-module}
Let $M$ be the orbit-closed Boolean/Clifford feature module generated by all four
feature types over the $15$ nonzero Boolean masks.  Then $M$ has $60$ feature
columns and is closed under coordinate permutations.  Projection by the $W(3,3)$
negative spectral projector gives
\[
  \operatorname{rank}(P_-M)=15.
\]
Therefore the $W(3,3)$ negative eigenspace is realized as the $15$-dimensional
image, equivalently a quotient, of the orbit-closed Boolean/Clifford feature
module.  This is the corrected form of the bridge: the selected BT1155 columns
need not be permuted literally, but the orbit-closed module has an invariant
projected image equal to the negative sector.
\end{theorem}
